\section[title={2024-08-27 - Postmodernes
Erzählen},reference={2024-08-27-Postmodernes Erzählen}]

\startframedtask
Lesen Sie den Ausschnitt und {\bf fassen} Sie die wesentlichen {\bf
Merkmale} postmodernen Erzählens stichwortartig {\bf zusammen}

\stopframedtask

\startitemize[packed]
\item
  große Komplexität (sprachlich + figural)
\item
  Kein ausgeprägter Charakter
\item
  ironischdistanziert
\item
  liebt das Parodistische
\item
  weltanschaulich nicht festgelegt
\item
  Unterhaltungsliteratur
  \startitemize[packed]
  \item
    Reizvolle Themen
  \stopitemize
\item
  Postmoderne: Ab 1990 bis jetzt
\stopitemize

\startframedtask
{\bf Beurteilen Sie} ob \quotation{Das Parfum} als postmoderner Roman
dargestellt wird.

\stopframedtask

\startitemize[packed]
\item
  Historischer / Exotischer Handlungsort
\item
  \quotation{Verbotene} Themen werden angesprochen
  \startitemize[packed]
  \item
    Mord
  \item
    Sex, Lust
  \stopitemize
\item
  Auflösung der Genre grenzen
  \startitemize[packed]
  \item
    Roman / Thrilla
  \stopitemize
\item
  Keine moralisierende Botschaft
  \startitemize[packed]
  \item
    Keine Kritik an seinem Tun (den Morden)
  \stopitemize
\stopitemize

\startframedtask
{\bf Stellen} Sie die wesentlichen Merkmale der verschiedenen Zäsuren im
Hinblick auf die Literatur und Sprache {\bf dar} und grenzen diese
voneinender ab.

\stopframedtask

Trümmerliteratur markiert den Bruch mit der nationalsozialistischen
Ideologie und die Hinwendung zu einer realistischen Darstellung der
unmittelbaren Nachkriegsrealität.
